(1) 全链各段具有相同速率 $\dot x$，动能为 $\frac12\lambda L\dot x^2$。以桌面为零势能，只有悬垂段具有位能 $U(x)=-\lambda g x^2/2$。在无耗散且忽略折角储能的理想条件下，拉氏方程得 $\lambda L\ddot x=\lambda g x$，故 $\ddot x=(g/L)x$；有效惯性始终是整条链而不是仅悬垂段。
(2) 设 $\omega=\sqrt{g/L}$，初值 $x(0)=x_0$、$\dot x(0)=0$ 对应 $x(t)=x_0\cosh(\omega t)$、$\dot x=x_0\omega\sinh(\omega t)$。到达 $x=L$ 的第一时间为 $t_*= \sqrt{L/g}\,\operatorname{arcosh}(L/x_0)$，仅在 $0\le t\le t_*$ 时符合指定几何约束。
(3) 从初态能量 $-\lambda gx_0^2/2$ 与任意时刻 $E=\frac12\lambda L\dot x^2-\frac12\lambda g x^2$ 相等，得 $\dot x^2=(g/L)(x^2-x_0^2)$。将双曲函数解代入，右端为 $(g/L)x_0^2\sinh^2(\omega t)$，恰与速度平方相符；能量和动力学推导相互独立。
(4) $x_0\to0^+$ 时 $t_*\sim\sqrt{L/g}\ln(2L/x_0)\to\infty$，说明理想零悬垂长没有初始驱动力。$x_0\to L^-$ 时 $t_*\to0$。末端越过桌角以后自由端、弯曲波、碰撞或非理想耗散会改变运动，故不能把本解延拓至 $x>L$。
